% Frittstående fil - kompilerer på egen hånd, er IKKE input i IMAX3012.tex.
%
% "Materiale"-kolonnen inneholder en klikkbar lenke ("Komp.") til
% kompendium-PDF-en på GitHub, etterfulgt av vanlig tekst med kapittel-/
% delkapittelnummer, samt referanser til "Adam and Essex" (uthevet skrift
% markerer sentrale tema, jf. tabellhode). "Relevante oppgaver"-kolonnen
% lister oppgavenummer per (del)kapittel, ikke spesielt anbefalt - bare til
% orientering. Denne filen har ingen avhengighet til IMAX3012.tex/.aux og
% kan bygges helt uavhengig.
%
% Dokumentet er i LANDSKAPSFORMAT for å få plass til den ekstra
% oppgave-kolonnen og de utvidede materiale-listene.
%
% Modul-etiketten i venstre kolonne er sentrert vertikalt over hele
% modulens radspenn ved hjelp av usynlige TikZ-merker (\tikzmark) plassert
% i første og siste rad, og en TikZ-node tegnet midt mellom dem rett etter
% modulens siste rad (via \noalign, altså mellom rader, ikke inni en celle).
% Dette krever TO pdflatex-kjøringer for at merkene skal bli riktige
% (TikZ "remember picture"-mekanismen, samme prinsipp som \ref/\label).

\documentclass[10pt,a4paper]{article}
\usepackage[a4paper,landscape,margin=0.8cm]{geometry}
\usepackage[norsk]{babel}
\usepackage[utf8]{inputenc}
\usepackage[T1]{fontenc}
\usepackage{lmodern}
\usepackage{array}
\usepackage{xcolor}
\usepackage{colortbl}
\usepackage{longtable}
\usepackage{multirow}
\usepackage{enumitem}
\usepackage{hyperref}
\usepackage{tikz}
\usetikzlibrary{tikzmark,calc}

\definecolor{fpHeader}{HTML}{8FC1A8}
\definecolor{fpModul1}{HTML}{FDF6DC}
\definecolor{fpModul2}{HTML}{F1E4E1}
\definecolor{fpModul3}{HTML}{D2EEDC}
\definecolor{fpRep}{HTML}{F2F2F2}
\definecolor{fpMatBlue}{HTML}{3B6FA0}
\definecolor{fpBorder}{HTML}{9FAFA8}

\hypersetup{
  colorlinks = true,
  linkcolor = fpMatBlue,
  urlcolor = fpMatBlue,
  citecolor = fpMatBlue,
}

% Lenke til kompendium-PDF-en på GitHub. Brukes som "Komp." i hver celle i
% Materiale-kolonnen.
\newcommand{\fpKompURL}{https://github.com/imfdrift/IMAx3012/blob/main/Kompendium/IMAX3012.pdf}
\newcommand{\fpKomp}{\href{\fpKompURL}{Komp.}}

% Liten hjelpemakro for "Adam and Essex"-lista i Materiale-kolonnen, og for
% en sentrert Komp.-lenkelinje øverst i cellen.
\newcommand{\fpKompLine}[1]{{\centering\fpKomp\ #1\par}}
\newcommand{\fpAE}{\textbf{Adam and Essex:}}
\newenvironment{fpItems}
  {\begin{itemize}[nosep,leftmargin=1em,itemsep=1pt,topsep=2pt,parsep=0pt]}
  {\end{itemize}}
% Rene "linjer" uten punktmerke -- brukes i stedet for rå \\ inne i
% tabellceller, siden en bar \\ inni en m{}-kolonne kan i praksis bli
% tolket som slutten på hele tabellraden (ikke bare en linjeskift i
% cellen). itemize/list-miljøer er derimot trygge.
\newenvironment{fpLines}
  {\begin{itemize}[nosep,leftmargin=0pt,itemindent=0pt,label={},itemsep=1pt,topsep=1pt,parsep=0pt]}
  {\end{itemize}}

% Modul-etikett sentrert mellom to \tikzmark-punkter (se forklaring øverst).
\newcommand{\fpModulLabel}[3]{%
  \noalign{\begin{tikzpicture}[remember picture,overlay]
    \node[align=center,text width=2.0cm] at ($(pic cs:#1)!0.5!(pic cs:#2)$) {\footnotesize\textbf{#3}};
  \end{tikzpicture}}%
}

\pagestyle{empty}

\begin{document}

\begin{center}
{\Large\bfseries Fremdriftsplan}\\[.3em]
{\large IMAA3012 -- Matematikk 3B}
\end{center}
\vspace{.6em}

Tabellen under viser OMTRENTLIG fremdriftsplanen for kurset, uke for uke, med tilhørende tema, arbeidskrav og relevant materiale fra kompendiet (uthevet skrift i Materiale-kolonnen markerer sentrale tema i pensum fra Adam and Essex). Materialet ligger på GitHub og blir kontinuerlig oppdatert. Spesielt i modul 3 vil der forekomme en del endringer. Oppgavene i siste kolonne er ikke spesielt anbefalt, men ment som orientering.

\begingroup
\scriptsize
\setlength{\tabcolsep}{3pt}
\renewcommand{\arraystretch}{1.1}
\arrayrulecolor{fpBorder}

\begin{longtable}{
  |>{\centering\arraybackslash}m{2.2cm}
  |>{\centering\arraybackslash}m{0.8cm}
  |>{\raggedright\arraybackslash}m{4.6cm}
  |>{\centering\arraybackslash}m{1.8cm}
  |>{\raggedright\arraybackslash}m{9.2cm}
  |>{\raggedright\arraybackslash}m{7.0cm}|
}
\hline
\rowcolor{fpHeader}
\textbf{Modul} & \textbf{Uke} & \textbf{Tema} & \textbf{Relevant arbeidskrav} & \textbf{Materiale} \newline \footnotesize(uthevet skrift behandler sentrale tema) & \textbf{Relevante oppgaver} \newline \footnotesize(Ikke spesielt anbefalt) \\
\hline
\endfirsthead
\hline
\rowcolor{fpHeader}
\textbf{Modul} & \textbf{Uke} & \textbf{Tema} & \textbf{Relevant arbeidskrav} & \textbf{Materiale} \newline \footnotesize(uthevet skrift behandler sentrale tema) & \textbf{Relevante oppgaver} \newline \footnotesize(Ikke spesielt anbefalt) \\
\hline
\endhead

% ===================== MODUL 1: Multippel-integrasjon =====================
\cellcolor{fpModul1}\tikzmark{mod1top}
 & \cellcolor{fpModul1}34 &
 \cellcolor{fpModul1}
 \begin{itemize}[nosep,leftmargin=1em]
   \item \textbf{Forelesning 1}
   \begin{itemize}[nosep,leftmargin=1em]
     \item Introduksjon og repetisjon.
     \item Python som verktøy
   \end{itemize}
   \item \textbf{Forelesning 2}
   \begin{itemize}[nosep,leftmargin=1em]
     \item Dobbelintegral i kartesiske koordinater
     \item Iterasjon som metode
     \item Anvendelser
   \end{itemize}
 \end{itemize}
 & \cellcolor{fpModul1}\multirow{2}{1.8cm}{\centering Øving 1} & \cellcolor{fpModul1}
 \fpKompLine{kapittel 0 og 1.1}
 \vspace{2pt}\fpAE
 \begin{fpItems}
   \item Derivasjon: kap 2, særlig \textbf{2.1} om tangent / normal.
   \item Nivåkurver: \textbf{13.1}
   \item Integrasjon: Kap 5,6.
   \item omdreiningslegemer: 7.1-\textbf{7.3}.
   \item Dobbelint.: 15.1,\textbf{15.2}, (15.3)
   \item Lin.alg: 10.7, \textbf{determinanter}
   \item Partiellderivasjon: 13.2,13.4
   \item Tangentplan: 10.4, \textbf{13.3}
   \item Sylinder-/kulekoordinater: \textbf{10.6}
   \item Kjeglesnitt: 8.1
 \end{fpItems}
 & \cellcolor{fpModul1}
 \begin{fpLines}
   \item \textbf{10.1:} 3, 15, 19, 21, 29
   \item \textbf{10.3:} 1, 3, 5
   \item \textbf{10.4:} 3
   \item \textbf{10.5:} Gjenkjen og skisser figurene fra 1, 5, 7, 11, 13, 15, (17)
   \item \textbf{10.6:} 1, 3, 5, 7, 11
   \item \textbf{8.1:} 7, 9, 11, 13
 \end{fpLines}
 \vspace{4pt}
 \begin{fpLines}
   \item \textbf{13.1:} 1, 11, 15, 19, 21, 27, 37, 39
   \item \textbf{13.2:} 3, 7
   \item \textbf{13.3:} 3, 13
 \end{fpLines}
 \\
\cline{2-3}\cline{5-6}

\cellcolor{fpModul1} & \cellcolor{fpModul1}35 &
 \cellcolor{fpModul1}
 \begin{itemize}[nosep,leftmargin=1em]
   \item \textbf{Forelesning 3}
   \begin{itemize}[nosep,leftmargin=1em]
     \item Dobbeltintegraler i polarkoordinater
   \end{itemize}
   \item \textbf{Forelesning 4}
   \begin{itemize}[nosep,leftmargin=1em]
     \item Python
   \end{itemize}
 \end{itemize}
 & \cellcolor{fpModul1} & \cellcolor{fpModul1}
 \fpKompLine{kap.\ 1.2-1.3}
 \vspace{2pt}\fpAE
 \begin{fpItems}
   \item Polarkoord.: 8.5-8.6
   \item Dobbelint i polar: \textbf{15.4} (til side 855)
 \end{fpItems}
 & \cellcolor{fpModul1}
 \begin{fpLines}
   \item \textbf{8.5:} 3, 9, 13, 27
   \item \textbf{8.6:} 13
 \end{fpLines}
 \vspace{4pt}
 \begin{fpLines}
   \item \textbf{15.4:} 1, 5, 11, 19, 21
 \end{fpLines}
 \\
\cline{2-6}

\cellcolor{fpModul1} & \cellcolor{fpModul1}36 &
 \cellcolor{fpModul1}
 \begin{itemize}[nosep,leftmargin=1em]
   \item \textbf{Forelesning 5}
   \begin{itemize}[nosep,leftmargin=1em]
     \item Trippelintegral i kartesiske koordinater
     \item Iterasjon som metode
   \end{itemize}
   \item \textbf{Forelesning 6}
   \begin{itemize}[nosep,leftmargin=1em]
     \item Trippelintegraler i sylinderkoordinater
     \item Trippelintegraler i kulekoordinater
   \end{itemize}
 \end{itemize}
 & \cellcolor{fpModul1}\multirow{2}{1.8cm}{\centering Øving 2} & \cellcolor{fpModul1}
 \fpKompLine{kap.\ 1.4-1.5}
 \vspace{2pt}\fpAE
 \begin{fpItems}
   \item Trippelint.: 15.5-15.6 (fra nederst side 868 og utover)
 \end{fpItems}
 & \cellcolor{fpModul1}
 \begin{fpLines}
   \item \textbf{15.5:} 1, 3, 15, 17, 19
   \item \textbf{15.6:} 1, 3, 5
 \end{fpLines}
 \\
\cline{2-3}\cline{5-6}

\cellcolor{fpModul1} & \cellcolor{fpModul1}37 &
 \cellcolor{fpModul1}
 \begin{itemize}[nosep,leftmargin=1em]
   \item \textbf{Forelesning 7}
   \begin{itemize}[nosep,leftmargin=1em]
     \item Variabelbytte generelt (Jacobi)
   \end{itemize}
   \item \textbf{Forelesning 8}
   \begin{itemize}[nosep,leftmargin=1em]
     \item Python, / anvendelser / oppsamling
   \end{itemize}
 \end{itemize}
 & \cellcolor{fpModul1} & \cellcolor{fpModul1}
 \fpKompLine{kap.\ 1.5-1.6}
 \vspace{2pt}\fpAE
 \begin{fpItems}
   \item Variabelbytte: \textbf{15.4} (s.855-860), 15.6 (867-868).
   \item Anvendelser: Kap 15.7 (876-882)
 \end{fpItems}
 & \cellcolor{fpModul1}
 \begin{fpLines}
   \item \textbf{15.4:} 1, 5, 11, 19, 21
 \end{fpLines}
 \vspace{4pt}
 \begin{fpLines}
   \item \textbf{15.7:} 19
 \end{fpLines}
 \\
\cline{2-6}

\cellcolor{fpModul1}\tikzmark{mod1bot} & \cellcolor{fpModul1}38 &
 \cellcolor{fpModul1}
 \begin{itemize}[nosep,leftmargin=1em]
   \item \textbf{Forelesning 9}
   \begin{itemize}[nosep,leftmargin=1em]
     \item Parametrisering av kurver
     \item vektor-valuerte funksjoner
     \item buelengde-integralet
     \item Linjeintegral for funksjoner
   \end{itemize}
   \item \textbf{Forelesning 10}
   \begin{itemize}[nosep,leftmargin=1em]
     \item parameterflater
     \item Flateintegral
   \end{itemize}
 \end{itemize}
 & \cellcolor{fpModul1} & \cellcolor{fpModul1}
 \fpKompLine{kap.\ 1.7-1.8}
 \vspace{2pt}\fpAE
 \begin{fpItems}
   \item Buelengde: 7.3
   \item Parametrisering av kurver: 8.2,8.3,8.4,12.1, \textbf{12.3}
   \item Linjeintegral: \textbf{16.3}
   \item Flater: 15.7,\textbf{16.5}
 \end{fpItems}
 & \cellcolor{fpModul1}
 \begin{fpLines}
   \item \textbf{8.2:} 1, 5, 7, (15)
   \item \textbf{8.3:} 3, 9, 17
   \item \textbf{8.4:} 1, (5)
 \end{fpLines}
 \vspace{4pt}
 \begin{fpLines}
   \item \textbf{12.1:} 3, 7, (15)
   \item \textbf{12.3:} 1, 7, 9, 13
 \end{fpLines}
 \vspace{4pt}
 \begin{fpLines}
   \item \textbf{15.7:} 1
 \end{fpLines}
 \vspace{4pt}
 \begin{fpLines}
   \item \textbf{16.3:} 1, 11
   \item \textbf{16.5:} 3, 7
 \end{fpLines}
 \\
\hline
\fpModulLabel{mod1top}{mod1bot}{Modul 1\\[2pt]\textit{Multippel-integrasjon}}

% ===================== MODUL 2: Vektoranalyse =====================
\cellcolor{fpModul2}\tikzmark{mod2top}
 & \cellcolor{fpModul2}39 &
 \cellcolor{fpModul2}
 \begin{itemize}[nosep,leftmargin=1em]
   \item \textbf{Forelesning 11}
   \begin{itemize}[nosep,leftmargin=1em]
     \item Vektorfelt
     \item Gradient, divergens og rotasjon
   \end{itemize}
   \item \textbf{Forelesning 12}
   \begin{itemize}[nosep,leftmargin=1em]
     \item Konservative vektorfelt
     \item Potensialfunksjoner
   \end{itemize}
 \end{itemize}
 & \cellcolor{fpModul2}Øving 3 & \cellcolor{fpModul2}
 \fpKompLine{kap 2.1, 2.2 og 2.4}
 \vspace{2pt}\fpAE
 \begin{fpItems}
   \item Vektorfelt: 16.1
   \item Gradient, div. curl: 13.7,\textbf{17.1},17.2
   \item Konserv. felt: \textbf{16.2}
 \end{fpItems}
 & \cellcolor{fpModul2}
 \begin{fpLines}
   \item \textbf{13.7:} 1, 3, (5,) 7, (9,) 11
 \end{fpLines}
 \vspace{4pt}
 \begin{fpLines}
   \item \textbf{16.1:} 1, 3
   \item \textbf{16.2:} 1, 3, 5
 \end{fpLines}
 \vspace{4pt}
 \begin{fpLines}
   \item \textbf{17.1:} 1, 5, 7
   \item \textbf{17.2:} (17)
 \end{fpLines}
 \\
\cline{2-6}

\cellcolor{fpModul2} & \cellcolor{fpModul2}40 &
 \cellcolor{fpModul2}
 \begin{itemize}[nosep,leftmargin=1em]
   \item \textbf{Forelesning 13}
   \begin{itemize}[nosep,leftmargin=1em]
     \item Linjeintegraler i vektorfelt
     \item Arbeid og sirkulasjon
   \end{itemize}
   \item \textbf{Forelesning 14}
   \begin{itemize}[nosep,leftmargin=1em]
     \item Fluksintegral i 2D og 3D
   \end{itemize}
 \end{itemize}
 & \cellcolor{fpModul2}\multirow{2}{1.8cm}{\centering Øving 4} & \cellcolor{fpModul2}
 \fpKompLine{kap 2.3, 2.5, 2.7}
 \vspace{2pt}\fpAE
 \begin{fpItems}
   \item Linjeint. av vektorfelt: \textbf{16.4}
   \item Fluksintegral: \textbf{16.6}
 \end{fpItems}
 & \cellcolor{fpModul2}
 \begin{fpLines}
   \item \textbf{16.4:} 1, 3
   \item \textbf{16.6:} 1, 3, 5
 \end{fpLines}
 \\
\cline{2-3}\cline{5-6}

\cellcolor{fpModul2} & \cellcolor{fpModul2}41 &
 \cellcolor{fpModul2}
 \begin{itemize}[nosep,leftmargin=1em]
   \item \textbf{Forelesning 15}
   \begin{itemize}[nosep,leftmargin=1em]
     \item Greens teorem
   \end{itemize}
   \item \textbf{Forelesning 16}
   \begin{itemize}[nosep,leftmargin=1em]
     \item Gauss divergenssetning
   \end{itemize}
 \end{itemize}
 & \cellcolor{fpModul2} & \cellcolor{fpModul2}
 \fpKompLine{kap 2.6, 2.8}
 \vspace{2pt}\fpAE
 \begin{fpItems}
   \item Greens: \textbf{17.3}
   \item Gauss: \textbf{17.4}
 \end{fpItems}
 & \cellcolor{fpModul2}
 \begin{fpLines}
   \item \textbf{17.3:} 1, 3, 5
   \item \textbf{17.4:} 1, 3, (11, 13)
 \end{fpLines}
 \\
\cline{2-6}

\cellcolor{fpModul2}\tikzmark{mod2bot} & \cellcolor{fpModul2}42 &
 \cellcolor{fpModul2}
 \begin{itemize}[nosep,leftmargin=1em]
   \item \textbf{Forelesning 17}
   \begin{itemize}[nosep,leftmargin=1em]
     \item Stokes teorem
   \end{itemize}
   \item \textbf{Forelesning 18}
   \begin{itemize}[nosep,leftmargin=1em]
     \item Repetisjon og oppgaveregning med de store teoremene
   \end{itemize}
 \end{itemize}
 & \cellcolor{fpModul2} & \cellcolor{fpModul2}
 \fpKompLine{kap 2.9}
 \vspace{2pt}\fpAE
 \begin{fpItems}
   \item Stokes: \textbf{17.5}
 \end{fpItems}
 & \cellcolor{fpModul2}
 \begin{fpLines}
   \item \textbf{17.5:} 1, 7
 \end{fpLines}
 \\
\hline
\fpModulLabel{mod2top}{mod2bot}{Modul 2\\[2pt]\textit{Vektoranalyse}}

% ===================== MODUL 3: PDEer =====================
\cellcolor{fpModul3}\tikzmark{mod3top}
 & \cellcolor{fpModul3}43 &
 \cellcolor{fpModul3}
 \begin{itemize}[nosep,leftmargin=1em]
   \item \textbf{Forelesning 19}
   \begin{itemize}[nosep,leftmargin=1em]
     \item Introduksjon til PDEer
   \end{itemize}
   \item \textbf{Forelesning 20}
   \begin{itemize}[nosep,leftmargin=1em]
     \item Bevaringslover
   \end{itemize}
 \end{itemize}
 & \cellcolor{fpModul3}Øving 5 & \cellcolor{fpModul3}
 \fpKompLine{kap 3.1-3.2}
 \vspace{2pt}{\centering YouTube: Introduksjon\par}
 & \cellcolor{fpModul3}
 \\
\cline{2-6}

\cellcolor{fpModul3} & \cellcolor{fpModul3}44 &
 \cellcolor{fpModul3}
 \begin{itemize}[nosep,leftmargin=1em]
   \item \textbf{Forelesning 21}
   \begin{itemize}[nosep,leftmargin=1em]
     \item Numerikk for PDEer: Finite differanser og randverdiproblem i 1D
   \end{itemize}
   \item \textbf{Forelesning 22}
   \begin{itemize}[nosep,leftmargin=1em]
     \item Numerikk for PDEer: 2D
   \end{itemize}
 \end{itemize}
 & \cellcolor{fpModul3}\multirow{2}{1.8cm}{\centering Øving 6} & \cellcolor{fpModul3}
 \fpKompLine{kap 3.3}
 & \cellcolor{fpModul3}
 \\
\cline{2-3}\cline{5-6}

\cellcolor{fpModul3} & \cellcolor{fpModul3}45 &
 \cellcolor{fpModul3}
 \begin{itemize}[nosep,leftmargin=1em]
   \item \textbf{Forelesning 23}
   \begin{itemize}[nosep,leftmargin=1em]
     \item Endelig-volum-metode
   \end{itemize}
   \item \textbf{Forelesning 24}
   \begin{itemize}[nosep,leftmargin=1em]
     \item Endelig-volum-metode
   \end{itemize}
 \end{itemize}
 & \cellcolor{fpModul3} & \cellcolor{fpModul3}
 & \cellcolor{fpModul3}
 \\
\cline{2-6}

\cellcolor{fpModul3}\tikzmark{mod3bot} & \cellcolor{fpModul3}46 &
 \cellcolor{fpModul3}
 \begin{itemize}[nosep,leftmargin=1em]
   \item \textbf{Forelesning 25}
   \begin{itemize}[nosep,leftmargin=1em]
     \item Endelig-volum-metode
   \end{itemize}
   \item \textbf{Forelesning 26}
   \begin{itemize}[nosep,leftmargin=1em]
     \item Endelig-volum-metode
   \end{itemize}
 \end{itemize}
 & \cellcolor{fpModul3}\multirow{2}{1.8cm}{\centering Øving 7} & \cellcolor{fpModul3}
 Presentasjon av master ved Robin T.\ Bye:
 \begin{fpItems}
   \item \textit{MSMECAUT\_2025-11-18.pdf}
   \item \textit{MSMECAUT\_Brochure.pdf}
 \end{fpItems}
 & \cellcolor{fpModul3}
 \\
\cline{1-3}\cline{5-6}
\fpModulLabel{mod3top}{mod3bot}{Modul 3\\[2pt]\textit{PDEer}}

% ===================== Repetisjon =====================
\rowcolor{fpRep}
\textit{Repetisjon} & 47 &
 \begin{itemize}[nosep,leftmargin=1em]
   \item Repetisjon
 \end{itemize}
 & & & \\
\hline

\end{longtable}
\endgroup

\end{document}
